% Chapter 4: Analysis and comments.

\newpage
\section{Analysis and Comments}
\label{sec:analysis}

The results show that the implementation completes the course tasks and that
two scale paths are available. This chapter asks two further questions: what do
the measured costs mean asymptotically, and how far is the current engine from
a more general search system?

\subsection{Time and Space Complexity}
\label{sec:complexity}

Table~\ref{tab:complexity} gives upper bounds for the main operations. $N$,
$T$, $P$, $D$, and $H$ denote the number of tokens, terms, postings,
documents, and matching documents; $k$ is the number of query terms, $R$ the
number of temporary runs, $B$ the block budget, and $C$ the cost of one
dictionary comparison. In phrase verification, $\mathrm{tf}_{\max}$ is the
largest frequency of a candidate term within one document.

\begin{table}[!htbp]
  \centering
  \caption{Time and space upper bounds for the main operations.}
  \label{tab:complexity}
  \small
  \begin{tabular}{l l l}
    \toprule
    Operation & Time & Space \\
    \midrule
    Index build (in memory) & $O(N+\sum \mathrm{df}\cdot C)$ & $O(T+P+N)$ \\
    Index build (block) & $O(N+RB\log B+TR)$ &
      $O(B+\text{postings of the current run})$ \\
    Word query & $O(\log T+\mathrm{df}+H\log H)$ & $O(D)$ \\
    Phrase query &
      $O(\mathrm{df}_{\min}\,\mathrm{tf}_{\max}\,k\log\mathrm{tf}_{\max}+H\log H)$ &
      $O(D+k)$ \\
    Boolean query & $O(\sum \mathrm{df}+H\log H)$ & $O(D)$ \\
    v2 encoding & $O(P+N)$ & $O(\text{one block})$ \\
    B+ tree lookup & $O(h\log F)$, plus $O(h)$ page reads & $O(h)$ page cache \\
    \bottomrule
  \end{tabular}
  \par\smallskip
  \footnotesize Note: $F$ is the B+ tree fanout and $h$ its height.
\end{table}

In-memory index construction spends most of its time on one dictionary access per
token, and its output keeps the term table, the postings, and the positions at
the same time. The three dictionaries produce indexes of the same size; only
the comparison count and the comparison constant change. Block construction uses
$B$ to bound the live terms at any moment, but it does not remove the cost: a
smaller $B$ means more runs, and sorting and merging become visible. The disk
B+ tree makes a different trade: a query reads the $O(h)$ pages on the
root-to-leaf path and no longer requires the full term table in memory.

On the query side, the main cost is scanning the relevant postings and sorting
the matching documents. A phrase query starts from the term with the smallest
document frequency, which shrinks the candidate set, but it still verifies $k$
positions in every candidate. The thresholds relate to complexity directly:
\code{max-df} and \code{min-df} remove terms before the posting scan and can
reduce later work, whereas \code{min-match} and \code{top-x} act after scoring;
\code{top-x} only reduces the output and does not save the scan or the sort. In
this implementation the thresholds are first a tool for changing query
semantics, not a uniform performance switch.

\subsection{Comparison with Alternative Data Structures and Algorithms}
\label{sec:dict-analysis}

Table~\ref{tab:dict-compare} compares the three dictionary implementations on
the same formal index.

\begin{table}[!htbp]
  \centering
  \caption{Build cost of three dictionary implementations (formal
  configuration, three builds each).}
  \label{tab:dict-compare}
  \footnotesize
  \renewcommand{\arraystretch}{1.0}
  \begin{tabular}{l r r r l}
    \toprule
    Dictionary & Comparisons & Rotations & Rehashes &
      Build time (in-process / external median) \\
    \midrule
    Splay tree (main) & 6\,392\,724 & 5\,850\,795 & 0 & 0.19\,s / 0.2610\,s \\
    Hash table & 688\,494 & 0 & 11 & 0.09\,s / \textbf{0.1603\,s} \\
    \code{std::map} (control) & 9\,435\,694 & 0 & 0 & 0.14\,s / 0.2084\,s \\
    \bottomrule
  \end{tabular}
  \par\smallskip
  \footnotesize Note: comparison, rotation, and rehash counts are in-process
  counters; external wall-clock time includes process start and file I/O, so it
  is comparable only in order of magnitude.
\end{table}

The hash table is fastest, \code{std::map} is second, and the splay tree is
slowest. The splay tree pays for 5.85 million rotations, but its in-order
traversal directly yields a sorted term table. The project therefore uses the
splay tree as the main dictionary, the hash table as the speed baseline, and
\code{std::map} as the balanced-tree control. The second format reduces the
index to 1.73 million bytes, about 46.82\% of v1.

The v1 and v2 formats represent a different trade-off. Version 1 stores fixed-
width integers and is simple to parse; version 2 delta-codes document IDs and
positions and uses byte-aligned variable-length codes, cutting the formal index
by 53.18\% with no change in query output. On the 500,000-document corpus,
however, the term table dominates and the overall ratio is only 66.55\%. The
right choice therefore depends on whether storage or build memory is the
binding constraint.

The in-memory tree and the disk B+ tree answer different questions. The splay
tree is appropriate when the whole term table fits in memory and the build
order is unknown; the B+ tree sidecar is appropriate when the table does not
fit and a cold query must avoid loading it. The sidecar is 237,568 bytes for
14,937 terms, and a cold lookup reads 0.225\% of the bytes of a full load
(114 page reads versus 29,819 cache hits during the verification pass). The
B+ tree is not a replacement for the in-memory dictionary; it is the boundary
mechanism for the next scale.

The stemmer comparison is also informative. The implementation follows the
original Porter algorithm \cite{porter1980} and is checked against NLTK's
\code{MARTIN\_EXTENSIONS} variant over all 25,705 forms. The original paper's
examples describe local rule effects, not necessarily the final stem; the
canonical C implementation and Snowball add fixes and extensions. The 41
over-stemmed forms found during testing came from reading the paper's example
column as final output. The corrected implementation matches the extended
variant, which is the appropriate target for a reproducible pipeline: the
choice of variant matters less than stating which one is used.

\subsection{Generalization: A Second Corpus}
\label{sec:rfc}
\label{sec:granularity}

The second corpus contains IETF RFC 1035 (DNS), 5321 (SMTP), and 7231 (HTTP
semantics). A generic importer converts the three documents to plain text; the
manifest and the text layer are the only interface. The same C++ binary,
without any code change, performs statistics, indexing, querying, and the
independent cross-check.

\begin{table}[!htbp]
  \centering
  \caption{End-to-end measurements on the RFC corpus.}
  \label{tab:rfc}
  \footnotesize
  \begin{tabularx}{\linewidth}{l X X}
    \toprule
    Metric & Raw configuration (no filter, no stemming) &
      Formal configuration (58 stop words + stemming) \\
    \midrule
    Documents & 3 & 3 \\
    Tokens & 75\,250 & 45\,370 \\
    Terms & 4\,328 & 2\,660 \\
    Postings & 7\,006 & 4\,594 \\
    Index bytes & 510\,181 & 308\,164 \\
    \code{brute_force.py} cross-check & all 4\,328 terms match &
      all 2\,660 terms match \\
    \bottomrule
  \end{tabularx}
  \par\smallskip
  \footnotesize Note: the RFC corpus has 9,946 lines and 81,848
  whitespace-separated words.
\end{table}

The RFC corpus is the clearest demonstration that there is no universal
stop-word list. Applying the same three-feature rule $(0.40, 0.001, 0.93)$
derives 58 stop words for the RFCs, overlapping the 39-word Shakespeare list in
only 26. The RFC 2119 keywords \code{must}, \code{may}, and \code{should} occur
in all three documents (df ratio 1), but their rate entropies are 0.809, 0.899,
and 0.885, below 0.93, so they are not mistaken for stop words. Conversely,
\code{server}, \code{page}, and \code{rfc} have rate entropy close to 1 and are
included in the derived list.

Queries behave as expected: \code{zone} matches all three documents,
\code{postmaster} matches the SMTP document and normalizes to \code{postmast},
and in \code{name server} the term \code{server} is filtered by the 58-word
list; the program explains this and queries only \code{name}. Conceptually, the
engine and the corpus-production stage are separate: a new corpus needs new
producer rules and stop-word thresholds, but not changes to the index or query
code.

Generalization has limits. The RFC plain-text import does not clean page
footers, so \code{page} from \code{[Page N]} enters the vocabulary and becomes
a stop word. This noise belongs to the producer side, not the engine side. A
new corpus still requires rechecking tokenization, document granularity, and
the stop-word distribution; ``the same binary runs'' does not mean ``no data
engineering is needed.''

\subsection{Limitations and Possible Improvements}
\label{sec:threshold-semantics}

Table~\ref{tab:limits} lists the main boundaries of the current implementation.
Phrase semantics and the ranking model have the largest effect: positions are
defined on the filtered token stream, so \code{my lord} degenerates to a query
for \code{lord}, and a phrase made entirely of stop words returns nothing.
Ranking uses only the matched-term count and the total term frequency, without
idf or length normalization, so long plays tend to occupy the top positions.

\begin{table}[!htbp]
  \centering
  \caption{System limitations, impact, and improvement directions.}
  \label{tab:limits}
  \footnotesize
  \begin{tabularx}{\linewidth}{l X X}
    \toprule
    Limitation & Impact & Improvement \\
    \midrule
    Phrases match on the filtered token stream &
      Phrases containing stop words degenerate; all-stop-word phrases return
      nothing &
      Keep stop-word positions, or build a second phrase index without stop-word
      removal \\
    ASCII-only tokenizer & Unicode corpora cannot be used directly &
      UTF-8 decoding and Unicode-category tokenization, followed by a full
      cross-check rerun \\
    Unclean plain-text import & Boilerplate such as RFC page numbers enters the
      vocabulary &
      Clean on the producer side, keeping the step auditable and optional \\
    Ranking uses only matched-term count and total tf & No inverse document
      frequency or length normalization &
      Implement TF-IDF/BM25 and redesign the relevance experiment \\
    Single-threaded build and query & The 201 threshold measurements take about
      4.05\,s serially &
      Parallel posting scans; parallel block construction \\
    Fixed document granularity (whole work) & Scene-level granularity is not
      part of the formal selection experiment &
      Parameterize granularity and compare scene-level with merged-short-poem
      stop-word distributions \\
    No incremental index & Updating one document requires a full rebuild &
      Segment-based indexing with periodic merging \\
    Single machine, no distribution & A real 400M-word deployment is still
      bounded by build time and disk &
      Shard by term range for build and query \\
    \bottomrule
  \end{tabularx}
\end{table}

The dictionary comparison should be stated honestly. Hashing is clearly faster
than the splay tree for pure construction: about 0.19\,s for the splay tree and
0.09\,s for the hash table in-process, but the splay tree's in-order traversal
directly yields a sorted term table and satisfies the course requirement for a
self-implemented data structure. All three implementations write byte-identical
indexes, which shows that the dictionary changes the construction process, not
the persisted result.

The threshold result most easily misunderstood is that ``tightening'' a
parameter can reduce the result set. In fact \code{--max-df} and
\code{--min-df} filter query terms first: this can make a single-term query
empty, or it can weaken a multi-term conjunction and increase the result.
\code{--min-match} changes the Boolean condition, and \code{--top-x} is the only
true result truncation. Phrase queries are unaffected by term-side thresholds,
which follows from the requirement that the whole phrase match.

\subsection{Bonus Analysis: Scalability and Boundaries}
\label{sec:bonus-analysis}

The bonus results separate cleanly into measured components and extrapolated
ones. On the query side, the B+ tree is measured from 100,000 to 50,000,000
keys at 15.10 bytes per key and height 4; the extrapolation to 400,000,000 keys
adds only file size and page count, not a new structural assumption. On the
build side, block construction is measured on 500,000 documents and 4,908,747
terms, with peak memory of 1.01\,GiB versus 2.00\,GiB for the in-memory build.
The measured peak for the full in-memory index implies about 160\,GiB at
400,000,000 distinct words if scaled linearly by term count, which is why the
in-memory path cannot answer the second axis.

The block build trades memory for merge time. Of the 168.16\,s total, 125.08\,s
(74.4\%) is merging, because the current merge scans all 4,855 run cursors to
find the smallest term, about $O(T\cdot R)$. A heap-based $k$-way merge would
reduce this to $O(T\log R)$ and is the first implementation improvement to
make. The temporary space is the other cost: 4,166,556 bytes on the
Shakespeare corpus and a much larger run set at 500,000 documents.

Compression and the disk dictionary optimize different resources. Version 2
cuts the formal index to 46.82\% of v1 but raises the 500k build peak from 1.01
to 1.46\,GiB; the B+ tree sidecar reduces cold-query I/O to 0.225\% of a full
load but adds a separate file and a page cache. None of these mechanisms is
universally better; each is appropriate when its constrained resource is the
one that binds.

The honest conclusion is therefore narrower than ``a 400,000,000-word system
runs.'' The key components scale: the dictionary is measured to 50,000,000
keys and structurally covers 400,000,000, and the block builder removes the
assumption that the whole term table must fit in memory. A complete large-scale
system still lacks direct SPIMI-to-B+ tree integration, distributed sharding,
and fault recovery. Those are the next engineering steps, and they should be
reported separately from the measured single-machine range.

\subsection{Conclusions}
\label{sec:conclusions}

The project implements the four required tasks on a 195-document Shakespeare
corpus. Word-frequency statistics produce 866,626 body tokens and 25,705
distinct terms; a three-feature rule derives 39 stop words with an overlap of
35 against the manual list; the stemmed inverted index contains 14,937 terms,
117,704 postings, and 556,866 positions; word, phrase, and Boolean queries
return document IDs, and the four thresholds are exercised over 201
measurements. The results are consistent with the specification in
Chapter~\ref{sec:design} and with the test cases in Chapter~\ref{sec:results}.

The most valuable outcome is not a single wall-clock number but the
reproducibility chain. An independent Python implementation rebuilds the
postings from the same text layer; three dictionaries write identical bytes;
block construction and in-memory construction agree; the second index format is
restored by an independent parser; and the disk dictionary matches a full load
term by term. Each of these constraints exposed implementation errors that a
``the output looks plausible'' check would have missed, including the dangling
\code{Node*&} in the splay insert, the temporary-reference bug in the index
reader, and the shifted offsets in the v2 format.

The engine is not a general search system. Its tokenizer is ASCII-only, its
phrase semantics inherit stop-word removal, its ranking has no idf or length
normalization, and it has no incremental or distributed mode. These are stated
as limitations in Table~\ref{tab:limits}, not hidden behind the measured
numbers. Within the single-machine, batch-built scope, however, the behavior is
specified, measured, and reproducible.

The bonus work supports a component-level claim rather than an end-to-end one.
The B+ tree dictionary and the block builder each address one resource limit,
and both are measured well beyond the formal corpus. Integrating them directly,
adding a heap merge, and sharding by term range are the concrete next steps
toward the 400,000,000-word target.
